\section{Triton row sum：row larger than block}

前面的 fused softmax 假设一整行能放进一个 program，本节故意放宽这个条件。Source coverage: \texttt{triton\_row\_sum\_example}；当 row larger than block 时，需要沿 columns 做 tiling，把多个 tile 的局部和累积到 registers，最后再执行一次 reduction，同时处理尾块 mask 与数值类型。

\begin{knowledgebox}{课堂提示：先问一行能否放进一个 block}
课程从 fused softmax 的前提反推限制：若整行能放进一个 block，Triton 可直接完成 block 内 reduction；若 4096 列而 block size 只有 1024，就必须把 row 拆成 4 个 tiles，让 threads 迭代 tiles、累积 partial sums，最后再对各 thread accumulator 做 reduction。这个容量问题决定了算法结构，而不是一个事后调参选项。
\end{knowledgebox}

\begin{knowledgebox}{讲义提醒：先写正确的分块不变量}
优化长 row 时先声明 accumulator 表示“已经处理过的所有 tiles 之和”，再移动 pointer。若循环边界、mask 或 dtype 不清楚，增加 warps 和 block size 只会更快地产生错误结果；Triton 调优必须从语义不变量开始。
\end{knowledgebox}


\begin{figure}[H]
\centering
\includegraphics[width=0.86\textwidth]{triton-row-sum.png}
\caption{triton row sum.png}
\figsource{CS336 Spring 2026 Lecture 6 官方可执行讲义。}
\end{figure}
\begin{knowledgebox}{读图：triton row sum.png}
这张图用于把抽象的 kernel 代码连接到硬件执行。读图时先看数据位于 HBM、shared memory 还是 registers，再看线程/blocks 如何映射到 SM。性能优化的关键是让数据在更靠近计算单元的位置被复用，减少慢速 HBM 往返。
\end{knowledgebox}

\begin{knowledgebox}{row 不适合一个 block 时怎么办}
如果一行太长，不能一次放进一个 block，就把行切成 tiles。每个 program 遍历多个 tiles，局部累加到 accumulator，最后再做 block 内 reduction。这是 tiling 思想的最小例子。
\end{knowledgebox}
\subsection{本章小结}
本节说明了一个 kernel optimization 的共同模式：先确认语义正确，再测时间，再看 profiler，最后用 block、shared memory、tiling、fusion 或 layout 调整降低数据移动。

